\section{完备性}
\label{sec:metric-completeness}

用迭代或截断逼近一个未知对象时，我们通常只能看到近似值本身：后面的项彼此越来越接近，极限却正是要找的东西，事先并不知道它在哪里。要判断这样的点列有没有极限，就需要一个只用点列内部距离表述的条件，这就是柯西列。本节先定义柯西列和完备空间，再讨论子空间何时完备，最后说明怎样补入缺失的极限，以及完备性带来的Baire纲定理。

\subsection{完备空间}
\label{sec:metric-complete-spaces}

收敛点列的后项不仅接近极限，也彼此接近。把后一性质单独拿出来，就得到不提及极限的条件。

\begin{definition}[柯西列]
\label{def:metric-cauchy}
度量空间$(X,d)$中的点列$(x_n)$若满足：对任意$\varepsilon>0$，存在$N$，使$m,n\geq N$时$d(x_m,x_n)<\varepsilon$，则称$(x_n)$为\term{柯西列}（Cauchy sequence）。
\end{definition}

\begin{lemma}[柯西列的基本性质]
\label{lem:metric-cauchy-basic}
在度量空间中：
\begin{enumerate}
\item 收敛点列是柯西列；
\item 柯西列是有界的；
\item 有收敛子列的柯西列本身收敛，并且收敛到该子列的极限。
\end{enumerate}
\end{lemma}

\begin{proof}
(1)设$x_n\to x$。对$\varepsilon>0$取$N$使$n\geq N$时$d(x_n,x)<\varepsilon/2$，则$m,n\geq N$时$d(x_m,x_n)\leq d(x_m,x)+d(x,x_n)<\varepsilon$。

(2)取$\varepsilon=1$对应的$N$，则$n\geq N$时$x_n\in B(x_N,1)$；前面只有有限项，按定理~\ref{thm:metric-convergent-bounded}证明中的办法放大半径即可。

(3)设$x_{n_k}\to x$。对$\varepsilon>0$，取$N$使$m,n\geq N$时$d(x_m,x_n)<\varepsilon/2$，再取$k$使$n_k\geq N$且$d(x_{n_k},x)<\varepsilon/2$。于是$n\geq N$时$d(x_n,x)\leq d(x_n,x_{n_k})+d(x_{n_k},x)<\varepsilon$。
\end{proof}

柯西列不一定收敛，关键看极限是否在当前空间中。用$\sqrt2$的十进制截断得到有理数列$q_1=1.4$，$q_2=1.41$，$q_3=1.414$，……。当$m\geq n$时$q_m$与$q_n$的前$n$位小数相同，所以$\lvert q_m-q_n\rvert<10^{-n}$，这是$\mathbb Q$中的柯西列。在$\mathbb R$中它趋于$\sqrt2$；若它在$\mathbb Q$中趋于某个有理数$q$，那么在$\mathbb R$中也趋于$q$，由极限唯一得$q=\sqrt2$，矛盾。所以这个柯西列在$\mathbb Q$中没有极限。

\begin{definition}[完备度量空间]
\label{def:metric-complete}
若度量空间$(X,d)$中的每个柯西列都在$X$中收敛，则称$(X,d)$是\term{完备度量空间}（complete metric space）。
\end{definition}

实轴是完备的：实数列收敛当且仅当它是柯西数列，这是确界原理的推论，可参考\BookChapterRef{chap:limits-continuity}{卷一“极限与连续”章}的柯西收敛准则。由此$(\mathbb R^n,\rho_\infty)$也完备：$\rho_\infty$下的柯西列在每个坐标上都是实柯西数列，各坐标分别收敛，而第~\ref{sec:metric-convergence}节说明逐坐标收敛就是$\rho_\infty$下的收敛。离散度量空间也完备：取$\varepsilon=1/2$，柯西列从某项起两两距离小于$1/2$，即从某项起恒定。

不完备的例子同样常见。$\mathbb Q$不完备，上面的截断列就是证据；区间$(0,1)$不完备，$1/(n+1)$是其中的柯西列，在$\mathbb R$中趋于$0$，但$0$不在区间内。这两个例子中的缺口性质不同：$\mathbb Q$缺少无理数，$(0,1)$只缺两个端点。完备性还依赖度量而不只依赖拓扑，第~\ref{chap:metric-space-maps}章会在实轴上构造一个与通常度量拓扑相同、却不完备的度量。

\subsection{闭子空间的完备性}
\label{sec:metric-closed-complete}

子空间是否完备，取决于它是否保留了环境中柯西列的极限。第~\ref{sec:metric-closure-sequences}小节的闭集判据正好描述这件事。

\begin{theorem}[闭子空间与完备子空间]
\label{thm:metric-closed-complete}
设$(X,d)$是度量空间，$A\subseteq X$非空。
\begin{enumerate}
\item 若$X$完备且$A$是闭集，则$(A,d_A)$完备；
\item 若$(A,d_A)$完备，则$A$是$X$的闭集，这里不要求$X$完备。
\end{enumerate}
因此在完备空间中，子空间完备当且仅当它是闭集。
\end{theorem}

\begin{proof}
(1)设$(a_n)$是$A$中的柯西列。子空间距离与环境距离相同，所以它也是$X$中的柯西列，由$X$完备收敛到某个$x\in X$。$A$是闭集，按序列判据$x\in A$，于是$(a_n)$在$A$中收敛。

(2)设$x\in\overline A$，按定理~\ref{thm:metric-closure-sequential}取$a_n\in A$使$a_n\to x$。由引理~\ref{lem:metric-cauchy-basic}，$(a_n)$是柯西列，它在完备的$A$中收敛到某个$a\in A$；在$X$中它既趋于$x$又趋于$a$，由极限唯一得$x=a\in A$。所以$\overline A=A$。
\end{proof}

在$\mathbb R$中，$[0,1]$与$\mathbb Z$都是闭集，因而完备；$\mathbb Q$不是闭集，相应地不完备。第(1)条的环境完备性不能省略。在$X=\mathbb Q$中，$A=\mathbb Q\cap[0,2]$是闭集，因为它是闭集$[0,2]$与$\mathbb Q$的交，但截断列$q_n$全在$A$中且在$A$中没有极限。同样，$A=(0,1/2]$在$X=(0,1)$中是闭集，却不完备。闭是相对于环境的性质，完备是空间自身的性质；只有当环境本身不缺极限时，前者才能保证后者。

\subsection{完备化}
\label{sec:metric-completion}

不完备的空间缺少一些极限。补入它们的思路是：既然缺失的极限只能通过柯西列来指认，就把柯西列本身当作新点；两个柯西列若彼此越来越近，它们指认的是同一个缺失点，应当识别。第~\ref{sec:metric-quotient}小节的零距离识别正好完成这一步。

设$(X,d)$是度量空间，$\mathcal C$是$X$中全体柯西列组成的集合。对$(x_n),(y_n)\in\mathcal C$，由三角不等式两次，
\begin{equation*}
\lvert d(x_n,y_n)-d(x_m,y_m)\rvert\leq d(x_n,x_m)+d(y_n,y_m),
\end{equation*}
所以$d(x_n,y_n)$是实柯西数列。实轴完备，于是可以定义
\begin{equation}
p\bigl((x_n),(y_n)\bigr)=\lim_{n\to\infty}d(x_n,y_n).
\label{eq:metric-completion-pseudometric}
\end{equation}
$p$的对称性和三角不等式由$d$逐项继承后取极限得到，$p((x_n),(x_n))=0$，所以$p$是$\mathcal C$上的伪度量。但它不是度量：$(1/n)$与$(0)$在$\mathbb R$中是不同的柯西列，$p$却为零。

按第~\ref{sec:metric-quotient}小节，把$p$为零的柯西列识别，得到度量空间$(\hat X,\hat d)$，其中$\hat X=\mathcal C/{\sim}$，$(x_n)\sim(y_n)$当且仅当$d(x_n,y_n)\to0$。对$x\in X$，令$\iota(x)$为常值点列$(x,x,\ldots)$的等价类。

\begin{theorem}[完备化]
\label{thm:metric-completion}
对任意度量空间$(X,d)$，上面构造的$(\hat X,\hat d)$是完备度量空间，$\iota:X\to\hat X$满足$\hat d(\iota(x),\iota(y))=d(x,y)$，并且$\iota(X)$在$\hat X$中稠密。
\end{theorem}

\begin{proof}
保持距离由定义直接得到：常值点列的逐项距离恒为$d(x,y)$。

稠密性：设$\xi\in\hat X$由柯西列$(x_n)$代表。对$\varepsilon>0$取$N$使$m,n\geq N$时$d(x_m,x_n)<\varepsilon$，则$m\geq N$时
\begin{equation*}
\hat d(\iota(x_m),\xi)=\lim_{n\to\infty}d(x_m,x_n)\leq\varepsilon.
\end{equation*}
所以$\iota(x_m)\to\xi$，$\xi$是$\iota(X)$中点列的极限。

完备性：设$(\xi_k)$是$\hat X$中的柯西列。由稠密性，对每个$k$取$x_k\in X$使$\hat d(\iota(x_k),\xi_k)<1/k$。由三角不等式与保持距离，
\begin{equation*}
d(x_k,x_l)=\hat d(\iota(x_k),\iota(x_l))<\tfrac1k+\hat d(\xi_k,\xi_l)+\tfrac1l,
\end{equation*}
右端在$k,l$充分大时任意小，所以$(x_k)$是$X$中的柯西列，记它的等价类为$\xi$。由稠密性证明中的估计，$\hat d(\iota(x_k),\xi)=\lim_l d(x_k,x_l)\to0$（$k\to\infty$），于是
\begin{equation*}
\hat d(\xi_k,\xi)\leq\hat d(\xi_k,\iota(x_k))+\hat d(\iota(x_k),\xi)<\tfrac1k+\hat d(\iota(x_k),\xi)\to0.
\end{equation*}
所以$\xi_k\to\xi$，$\hat X$完备。
\end{proof}

称$(\hat X,\hat d)$为$X$的\term{完备化}（completion）。保持距离的映射$\iota$让我们可以把$X$看成$\hat X$的子空间，$\hat X$中的新点恰好是$X$中柯西列的极限。若$X$本身完备，每个柯西列都等价于其极限的常值点列，$\iota$是满射，完备化不增加新点。

完备化在只差一个保持距离的双射的意义下唯一：若另有完备空间包含$X$的一个保持距离的稠密副本，它与$\hat X$之间存在保持距离的双射。证明需要把稠密子集上的映射延拓到全空间，留到第~\ref{chap:metric-space-maps}章。按这个唯一性，$\mathbb Q$的完备化就是$\mathbb R$，因为$\mathbb R$完备且$\mathbb Q$在其中稠密；$(0,1)$的完备化是$[0,1]$。需要注意，上面的构造在式\eqref{eq:metric-completion-pseudometric}中用到了实轴的完备性，所以它不能原样用来从$\mathbb Q$构造实数。

\subsection{Baire纲定理}
\label{sec:metric-baire}

两个稠密开集的交仍是稠密开集。设$U,V$稠密开，$W$是非空开集，则$W\cap U$是非空开集，再与稠密的$V$相交非空，所以$U\cap V$与每个非空开集相交。这一论证只用了拓扑，对有限个稠密开集都成立。可数个稠密开集的交却可能是空集。$\mathbb Q$中没有孤立点，所以对每个有理数$q_k$，$U_k=\mathbb Q\setminus\{q_k\}$都是$\mathbb Q$的稠密开集，而把有理数排成$q_1,q_2,\ldots$后，$\bigcap_k U_k=\varnothing$。完备性恰好排除了这种情形，这就是\term{Baire纲定理}（Baire category theorem）。

\begin{theorem}[Baire纲定理]
\label{thm:metric-baire}
设$(X,d)$是非空完备度量空间，$U_1,U_2,\ldots$是$X$中的稠密开集。则$\bigcap_{k\geq1}U_k$在$X$中稠密，特别地它非空。
\end{theorem}

\begin{proof}
任取非空开集$W$，要找一点属于$W\cap\bigcap_kU_k$。思路是选一列嵌套的闭球，每一步都落进下一个$U_k$，并让半径趋于零。

$W\cap U_1$是非空开集，取其中一点$x_1$及$r_1\in(0,1)$使$B(x_1,2r_1)\subseteq W\cap U_1$，于是闭球$B[x_1,r_1]\subseteq W\cap U_1$。设$x_k,r_k$已经选好。开球$B(x_k,r_k)$与稠密开集$U_{k+1}$的交是非空开集，取其中一点$x_{k+1}$及$r_{k+1}\in(0,r_k/2)$，使
\begin{equation*}
B[x_{k+1},r_{k+1}]\subseteq B(x_k,r_k)\cap U_{k+1}.
\end{equation*}
这样得到嵌套闭球$B[x_1,r_1]\supseteq B[x_2,r_2]\supseteq\cdots$，并且$r_k<2^{1-k}$。

当$m\geq k$时$x_m\in B[x_k,r_k]$，所以$d(x_m,x_k)\leq r_k\to0$，$(x_k)$是柯西列。由完备性，它收敛到某个$x\in X$。对固定的$k$，点列从第$k$项起都在闭集$B[x_k,r_k]$中，极限也在其中。于是$x\in B[x_k,r_k]\subseteq U_k$对每个$k$成立，并且$x\in B[x_1,r_1]\subseteq W$。
\end{proof}

证明中闭球的作用是让极限留在球内，半径减半的作用是让球心组成柯西列，完备性则保证这个柯西列有极限。$\mathbb Q$的反例失败在最后一步：嵌套球的球心可以逼近一个无理数。

取补集得到常用的等价形式。若闭集$F$的内部为空，则它的补集与每个非空开集相交，是稠密开集；反之亦然。称内部为空的闭集为\term{无处稠密}（nowhere dense）闭集，Baire纲定理等价于：在非空完备度量空间中，可数个无处稠密闭集之并的内部仍为空，特别地它们不能覆盖整个空间。

这个形式立刻给出一个熟悉的结论：$\mathbb R$不可数。若$\mathbb R=\{x_1,x_2,\ldots\}$，每个单点集$\{x_k\}$都是内部为空的闭集，可数个它们的并却是整个完备空间$\mathbb R$，与定理矛盾。对照$\mathbb Q$，它确实是可数个单点集之并，与$\mathbb Q$不完备相一致。

Baire纲定理的典型用法是从逐点的有界性得到某个球上的统一有界性。

\begin{theorem}[逐点有界蕴含局部一致有界]
\label{thm:metric-baire-pointwise-bound}
设$(X,d)$是非空完备度量空间，$\{f_i\}_{i\in I}$是一族连续函数$f_i:X\to\mathbb R$，并且对每个$x\in X$，$\sup_{i\in I}\lvert f_i(x)\rvert<\infty$。则存在开球$B(x_0,r)$和常数$M$，使对所有$i\in I$和$x\in B(x_0,r)$都有$\lvert f_i(x)\rvert\leq M$。
\end{theorem}

\begin{proof}
对正整数$k$，令$F_k=\{x:\lvert f_i(x)\rvert\leq k\text{ 对所有 }i\in I\}=\bigcap_{i\in I}f_i^{-1}([-k,k])$。连续映射下闭集的原像是闭集，任意交仍是闭集，所以$F_k$是闭集。由逐点有界，每个$x$属于某个$F_k$，即$X=\bigcup_kF_k$。按Baire纲定理，某个$F_k$的内部非空，它包含一个开球$B(x_0,r)$，取$M=k$即可。
\end{proof}

这里只用到度量与完备性，没有用到任何线性结构。对线性映射族，球上的有界性可以平移和伸缩到整个单位球，得到一致有界原理，这在第~\ref{chap:linear-operators}章展开。
